% Generative-models post, figure 1: discriminative vs generative models, after
% Shervine Amidi's CS229 cheatsheet table (recovered via the Wayback Machine,
% i.stack.imgur.com/Xrmqg.png). Same rows: goal, what is learned, an
% illustration, examples. Both illustrations use the same toy points (fa, fd):
% the discriminative model learns only the boundary, the generative model
% learns each class's distribution (nested density contours).
\documentclass[tikz,border=3pt]{standalone}
\input{FIGKIT/preamble.tex}
\begin{document}
\begin{tikzpicture}[fig,
    rowlab/.style={lab, anchor=base west, text=soft},
    cell/.style={anchor=base west, font=\rmfamily\small},
    head/.style={anchor=base west, font=\rmfamily\small\bfseries},
    hl/.style={draw=rule, line width=.5pt},
    pt/.style={circle, minimum size=2.6mm, inner sep=0pt, tint=#1, fill opacity=.4, line width=.5pt},
  ]
  % columns: labels at x=0, discriminative at 26mm, generative at 89mm; width 150mm
  \def\cA{26mm}\def\cB{89mm}
  \node[head] at (\cA,-1mm) {Discriminative model};
  \node[head] at (\cB,-1mm) {Generative model};
  \draw[hl] (0,-4mm) -- (150mm,-4mm);
  \node[rowlab] at (0,-10mm) {goal};
  \node[cell] at (\cA,-10mm) {Directly estimate $p(y\mid x)$};
  \node[cell] at (\cB,-10mm) {Estimate $p(x\mid y)$, then deduce $p(y\mid x)$};
  \draw[hl] (0,-14mm) -- (150mm,-14mm);
  \node[rowlab] at (0,-20mm) {learns};
  \node[cell] at (\cA,-20mm) {A decision boundary};
  \node[cell] at (\cB,-20mm) {The distribution of each class};
  \draw[hl] (0,-24mm) -- (150mm,-24mm);
  \node[rowlab] at (0,-44mm) {illustration};
  \draw[hl] (0,-62mm) -- (150mm,-62mm);
  \node[rowlab] at (0,-68mm) {examples};
  \node[cell] at (\cA,-68mm) {Regressions, SVMs};
  \node[cell] at (\cB,-68mm) {GDA, Naive Bayes};

  % the toy points, in a 56 x 32 mm panel; origin = panel centre
  \def\blue{-20.5/5.6, -16.2/7.1, -11/10.1, -5.4/10.1, -11.9/5.3, -6/5.3,
    -17.3/0.9, -13.1/0.9, -9.8/-0.6, -6/-0.6, -21.8/-3, -11/-4.7, -15/-6.6}
  \def\red{19.4/3.9, 15.4/2.3, 19.4/-0.6, 11/-0.9, 17.9/-2.7, 13.9/-5.9,
    18/-5.9, 2.7/-9.5, 6.2/-9.5, 12.7/-9.5, 16.5/-10.1}

  % discriminative: points + a boundary
  \begin{scope}[shift={(\cA+28mm,-43mm)}]
    \foreach \x/\y in \blue \node[pt=fa] at (\x mm,\y mm) {};
    \foreach \x/\y in \red  \node[pt=fd] at (\x mm,\y mm) {};
    \draw[draw=ink, line width=.8pt, dash pattern=on 3pt off 2.5pt] (-10mm,-15mm) -- (13mm,14.5mm);
  \end{scope}

  % generative: the same points + one density per class
  \begin{scope}[shift={(\cB+28mm,-43mm)}]
    \begin{scope}
      \clip[rounded corners=2.5pt] (-30mm,-17mm) rectangle (30mm,17mm);
      \foreach \s in {1, .68, .38} {
        \draw[tint=fa, fill opacity=.13, line width=.5pt, rotate around={-28:(-12.5mm,2.2mm)}]
          (-12.5mm,2.2mm) ellipse ({15.5mm*\s} and {11mm*\s});
        \draw[tint=fd, fill opacity=.13, line width=.5pt, rotate around={30:(13.8mm,-4.4mm)}]
          (13.8mm,-4.4mm) ellipse ({15mm*\s} and {9mm*\s});
      }
    \end{scope}
    \foreach \x/\y in \blue \node[pt=fa] at (\x mm,\y mm) {};
    \foreach \x/\y in \red  \node[pt=fd] at (\x mm,\y mm) {};
  \end{scope}
\end{tikzpicture}
\end{document}
